\chapter{Foundations: Learning Paradigms}
\label{chap:objectives}

\section{Types of Supervision}
\label{sec:objectives-supervised-unsupervised-self-supervised-and-reinforcement-learning}

\begin{table}[htbp]
    \centering
    \small
    \begin{tabularx}{\linewidth}{@{}L{.18\linewidth}YY@{}}
    \toprule
    Paradigm & Training information & Typical goal \\
    \midrule
    Supervised & Input--target pairs & Classification, regression, forecasting \\
    Unsupervised & Inputs without task labels & Density, clustering, compression \\
    Self-supervised & Targets manufactured from inputs & Representation learning from structure in the data \\
    Semi-supervised & Few labels plus many unlabelled examples & Label-efficient prediction \\
    Reinforcement learning & States, actions and rewards through interaction & Policy maximising expected return \\
    \bottomrule
    \end{tabularx}
    \caption{Learning paradigms mainly differ in where the training signal comes from.}
\end{table}

\section{Supervised Losses: Continuous}
\label{sec:objectives-supervised-losses-for-continuous-targets}
For residuals $e_i=y_i-\hat y_i$ over $N_{\text{examples}}$:

\paragraph{MSE} makes large residuals disproportionately expensive and matches constant-variance Gaussian observation noise up to scale/constants.
\[
    \mathrm{MSE}&=\frac{1}{N_{\text{examples}}}\sum_i e_i^2.
\]
\paragraph{MAE} grows linearly with residual size and matches a Laplace-like noise model.
\[
    \mathrm{MAE}&=\frac{1}{N_{\text{examples}}}\sum_i |e_i|.
\]
\paragraph{RMSE} restores target units and is often more useful for reporting than for changing the optimiser itself. 
See \cref{sec:evaluation-regression-metrics}.
\[
    \mathrm{RMSE}&=\sqrt{\mathrm{MSE}}.
\]
\paragraph{Huber} loss is quadratic near zero and linear for large residuals:
\[
    \ell_{\delta}(e)
    =
    \begin{cases}
    \frac12e^2,&|e|\le\delta,\\
    \delta\left(|e|-\frac12\delta\right),&|e|>\delta.
\end{cases}
\]
It preserves a smooth small-error region while reducing the influence of large outliers relative to MSE.

\begin{calloutmemoryhook}
MSE squares outliers and MAE does not. Huber is quadratic near zero and linear far away.
\end{calloutmemoryhook}

\section{Supervised Losses: Discrete}
\label{sec:objectives-classification-losses-and-logits}
A logit is a number between which we map to a probability, for example we might say that a score out of 100 of 34 is more likley than not a 0 in a binary system.
A binary logit $a$ is $p=\sigma(a)$ and,
\paragraph{Binary Cross-Entropy} is
\[
    \ell_{\mathrm{BCE}}
    =
    -\underbrace{y\log p}_{\text{reward positive probability when }y=1}
    -\underbrace{(1-y)\log(1-p)}_{\text{reward negative probability when }y=0}.
\]
Use implementations that combine the logit and BCE calculation for numerical stability.
Independent sigmoid/BCE outputs suit labels that can be simultaneously true.

For $N_{\text{classes}}$ mutually exclusive classes, softmax converts logits to a probability:
\[
    p_i
    =
    \frac{e^{a_i}}{\sum_{j=1}^{N_{\text{classes}}}e^{a_j}}.
\]
\paragraph{Categorical Cross-Entropy} for the true class is $-\log p_{\text{true class}}$.
Training loss and deployment metric are separate design choices; accuracy, precision/recall and F1 are in \cref{sec:evaluation-classification-evaluation-the-confusion-matrix,sec:evaluation-precision-recall-and-f1}.


\subsection{Differentiating Categories}
\label{sec:models-gumbel-softmax-differentiable-categorical-choices}
A hard categorical sample blocks ordinary pathwise gradients.
Gumbel--Softmax replaces the hard one-hot choice during optimisation with a temperature-controlled softmax of noisy category logits:
\[
    \underbrace{y_i}_{\text{soft category weight}}
    =
    \operatorname{softmax}_i\!\left(
    \frac{\underbrace{\log\pi_i}_{\text{category log-weight}}+\underbrace{g_i}_{\text{Gumbel noise}}}{\underbrace{\tau}_{\text{temperature}}}
    \right).
\]
Lower temperature makes the vector more nearly one-hot but also makes gradients more brittle.
A straight-through variant uses a hard forward choice with an approximate backward gradient.

\section{Supervised Losses: Ranking and Preference}
\label{sec:objectives-ranking-and-preference-objectives}
Ranking differs from ordinary regression because the useful output is an ordering rather than a numerically calibrated score, which exposes many different dimensions we may evaluate.
\begin{center}
    \small
    \begin{tabularx}{.96\linewidth}{@{}L{.18\linewidth}L{.28\linewidth}Y@{}}
    \toprule
    \textbf{Objective} & \textbf{Training question} & \textbf{Main trade-off} \\
    \midrule
    Pointwise & Is this item good? & Simple, but treats ranking as independent prediction \\
    Pairwise & Should item $a$ rank above $b$? & Direct ordering signal; creates many possible pairs \\
    Listwise & Is the whole ordered list useful? & Closest to ranking metrics; optimisation is less direct \\
    \bottomrule
    \end{tabularx}
\end{center}

A pairwise model needs only score differences:
\[
    \underbrace{P(a\succ b)}_{\text{$a$ should rank above $b$}}
    =
    \sigma\!\left(
    \underbrace{s_a-s_b}_{\text{relative score}}
    \right).
\]
This same relative-preference structure reappears in LLM preference optimisation in \cref{sec:llms-preference-optimisation}.

Listwise metrics such as \textbf{NDCG} care much more about mistakes near the top of the ranking.
Sorting is discrete, so small score changes can leave NDCG unchanged until two items swap order.
Therefore, methods such as \textbf{LambdaRank} use smooth pairwise updates whose strength reflects how much a swap would change the ranking metric.
The boosted-tree form \textbf{LambdaMART} is linked in \cref{subsec:trees-lambdamart}.

\section{Supervised Losses: Representation}
\label{sec:objectives-representation-learning-objectives}
\paragraph{Reconstruction} objective asks an encoder/decoder pair to preserve information needed to rebuild the input:
\[
    \mathcal L_{\text{reconstruction}}
    =
    \left\|x-g_{\theta_{\text{decoder}}}(f_{\theta_{\text{encoder}}}(x))\right\|_2^2.
\]
Reconstruction does not guarantee preservation of task-relevant information.

\paragraph{Contrastive} objectives instead shape relative geometry.
A triplet-style objective says ``positive closer than negative by a margin'':
\[
    \mathcal L_{\text{triplet}}
    =
    \max\left(0,
    \underbrace{d(z_{\text{anchor}},z_{\text{positive}})}_{\text{pull together}}
    -
    \underbrace{d(z_{\text{anchor}},z_{\text{negative}})}_{\text{push apart}}
    +
    \underbrace{m_{\text{margin}}}_{\text{required separation}}
    \right).
\]
The important experimental choices are how positives/negatives are constructed, which distance is used, and whether augmentations preserve task signal.

\section{Supervised Losses: Variational/Generative}
\label{sec:objectives-variational-and-generative-objectives}
A VAE balances data fit with a structured latent distribution:
\[
    \underbrace{\mathcal L_{\mathrm{VAE}}}_{\text{negative ELBO}}
    =
    \underbrace{-\E_{q_{\theta_{\text{encoder}}}(z\mid x)}\log p_{\theta_{\text{decoder}}}(x\mid z)}_{\text{reconstruction / likelihood}}
    +
    \underbrace{D_{\mathrm{KL}}\!\left(q_{\theta_{\text{encoder}}}(z\mid x)\|p(z)\right)}_{\text{latent prior compatibility}}.
\]
GANs, diffusion models and normalising flows use different objectives; ``generative model'' does not imply this ELBO.

\section{Self-Supervised losses: Reinforcement-Learning}
\label{sec:objectives-reinforcement-learning-objectives}
A policy aims to maximise expected discounted reward:
\[
    \underbrace{J(\theta)}_{\text{policy objective}}
    =
    \E_{\pi_\theta}
    \left[
    \sum_{t=0}^{N_{\text{horizon}}}
    \underbrace{\gamma^t}_{\text{discount}}
    \underbrace{r_t}_{\text{reward}}
    \right].
\]
The optimisation loss may combine policy-gradient, value-error and entropy terms.
Reward, training loss and evaluation metric therefore need not be identical.
Entropy itself is defined in \cref{sec:foundations-entropy}.